\documentclass[12pt,a4paper]{article}
\usepackage{comment}
\usepackage[hidelinks]{hyperref}
\usepackage{listings}
\usepackage{graphicx}
\usepackage{subcaption}
\usepackage{float}

\begin{comment}
    \url to show urls inside the document

    https://www.tablesgenerator.com/#
    We can create tables with this page so its easier to portray our examples
\end{comment}

\title{%
    \vspace*{-5mm}\Huge GW2-SRS EXTRACT \\
    \vspace*{2mm}\Large powered by \LaTeX}

\author{\vspace*{-5mm}\large Daniel Lopez: \textbf{Extraction algorithm}}


\begin{document}

    \maketitle
    \section*{1.0 EXTRACT}

    \section*{\large 1.1 Introduction}
    As the first part of the \textbf{ETL} it's important to have in mind our data sources. When
    I began with the project I only had one way of getting logs, and they were all coming
    from friends. This meant I had to manually parse them with EliteInsights\footnote{EliteInsights it's an app designed to parse \textit{.zevtc} files into JSON, HTML or CSV files.}. 

    This process was extremely slow, and when at first it was viable due to having few logs,
    it started being quite useless when the project meant to have lots of logs to analyze.
    The best option then was thinking on using web-scraping or an API connection; and thankfully
    we have a web-page available, thanks to \textit{Johannes Pfau}, where people can upload their logs.\\

    \noindent \href{https://gw2wingman.nevermindcreations.de/home}{GW2-Wingman} is probably the best approach I made to this project, not only in efficiency, but also
    in learn curve. Scraping was something I craved for since I started focusing on the Data Engineering
    route. It taught me a lot of things and problems I could solve, some easier, some harder, but
    all of them were at the end solved.

    \newpage

    \section*{\large 1.2 About GW2-Wingman architecture}
    As I was learning to scrape the web-page I realized someone wasn't right, and it made the scrape a
    bit harder when I understood the data was stored inside an \textbf{iframe tag}\footnote{An iframe is a nested web
    inside the main website}, so it was a matter of getting to the actual web-page.\\

    \noindent It was quite simple in the end, the iframe had an \textbf{href} tag containing the actual path, which only added
    \textit{/logContent} before the log name. \\

    \noindent This is a comparison of the actual URL showed, with the nested one:
    \begin{lstlisting}[language=Python]
    https://gw2wingman.nevermindcreations.de/log

    https://gw2wingman.nevermindcreations.de/logContent
    \end{lstlisting}

    \section*{\large 1.3 JSON data gathering}
    Extract algorithms are rather simple once you understand what you need to find, and that was
    the hardest part, since it required looking into the actual HTML code. I tackled the extraction
    with BeautifulSoup\footnote{As well as Selenium, BeautifulSoup let us extract information from any web-page.} 
    by passing an URL and the needed headers.\\

    \noindent After that I just needed to get the correct script and extract the JSON. This required a bit of 
    research; the JSON was nested on a JavaScript variable, which I had to strip off, and save only the
    JSON data.

    \begin{lstlisting}[language=Python]
    response = requests.get(url=url,headers=HEADERS)
    soup = BeautifulSoup(response.content,'html.parser')

    data = soup.find_all('script')[8]
    dataString = data.text.rstrip()
    \end{lstlisting}

    \newpage

    Last thing I needed to do is set up an algorithm, or in this case, a class, that help us know
    the boss Name and create a tag with it. This will make the further processes so much agile:

    \begin{lstlisting}[language=Python]
    class boss:

        def __init__(self, url: str) -> str:
            self.url = url

        def getBossName(self) -> str:
            url = self.url
            urlLines = url.split('/')
            if len(urlLines) < 5:
                bossName = urlLines[3]
            elif len(urlLines) == 5:
                bossName = urlLines[4]
            return bossName
        
        def getBossTag(self, bossName: str) -> str:
            bossTag = bossName.split('_')
            nameTag = bossTag[1]
            return nameTag
    \end{lstlisting}

    \section*{\large 1.4 URL copy code}
    Ending with the extraction part of the project, I will explain a quite relevant part here. 
    Normally we could extract information URL by URL, but it is not really convenient when we have, 
    like in this case, more than 100 logs. Therefore, best option is creating a script that copy and 
    write down this URLs for us on a text file we can later read.\\

    \begin{itemize}

        \item I set the search method to scope the a tags, so I only had to just take the href of each one and 
        add the url\_str and have the complete link for a later iteration.

        \item Nonetheless, the a tag contains about 5-6 lines that are not log hrefs, instead, they show an apikey 
        href and JavaScript void hrefs, so I just omited them and set a new line for each one the script finds.

        \newpage

        \item Finally, I wrote it down on a .txt file; however I realized that, at first I was using 'w' mode, but it 
        ended being quite inefficient since it constantly overwrites the file. To avoid this, I just simply changed 
        the 'w' mode to 'a'\footnote{An 'a' stands for append, and will add new data without overwriting the existing one}
        mode, this way we can just append the information and not worry about it overwriting over and over again.

    \end{itemize}

    This is how the code looks like:

    \begin{lstlisting}[language=Python]
    fh = open(path,'a')

    for link in soup.find_all('a'):
        url_str = 'https://gw2wingman.nevermindcreations.de'
        data = link.get('href')
        
        try:
            log_str = url_str+data
            if log_str.endswith('apikey'):
                log_str.replace('apikey','\n')
            elif log_str.endswith('void(0)'):
                log_str.replace('void(0)','\n')
            else:
                fh.write(log_str)
                fh.write('\n')
        except Exception as e:
            print('Error: ', str(e))
    \end{lstlisting}

    \begin{figure}[b]
        \centering
        \begin{subfigure}[b]{0.15\linewidth}
            \href{https://github.com/icharo-tb}{\includegraphics[width=\linewidth]{Images/GitHub-logo.png}}
        \end{subfigure}
        \begin{subfigure}[b]{0.15\linewidth}
            \href{https://www.linkedin.com/in/danielopezpajares/}{\includegraphics[width=\linewidth]{Images/LinkedIn_Logo.png}}
        \end{subfigure}
        \begin{subfigure}[b]{0.15\linewidth}
            \href{https://gitlab.com/daniel.lopez.pajares.2021}{\includegraphics[width=\linewidth]{Images/GitLab_logo.png}}
        \end{subfigure}
    \end{figure}
\end{document}